\section{Exact transfer to a bulk critical exponent}\label{sec:transfer}
Superimpose the finite-size engine on the ruler. Interior sites carry
the four tableau flags $T,F,B,R$, tied by two-site unit penalties to the
conditions that their upper, left, lower and right neighbours are the
corresponding $W,V,W,V$ walls. Wall sites carry one dummy tableau label.
Tableau rules act only between interior sites of one cell. Every
quantum site outside the cell's marked bottom-row segment is frozen.
The classical wall bonus replaces the physical-box flag reward inside
the cell; all additional tableau tests are nonnegative. In a perfect
lattice each cell consequently has the unique tableau of its
$h\times h$ box, where $h=p-1$.

Let $A$ be a fast density word from Lemma~\ref{lem:word}, modified so
$0,1,2\in A$. Let $D(h)$ denote the \emph{full} counter gap selected by
that tableau. With $\ell(h)=\bitlength(h-1)$ and
$k(h)=\beta|A\cap[0,\ell(h))|$, Lemma~\ref{lem:counter} gives
\begin{equation}\label{eq:cell-bracket}
 \frac{4^{-k(h)}}{64\ell(h)^2}\leq D(h)
 \leq4\pi^2 4^{-k(h)},\qquad
 D(h)\leq\frac{\pi^2}{\ell(h)^2}.
\end{equation}
The last inequality follows from the free orbit of length at least
$2\ell$: each summand in \eqref{eq:orbit} is at least one. It guarantees
gap closure even when the density, and hence the exponent, is zero.

\begin{proof}[Proof of Theorem~\ref{thm:bulk-headline}]
All classical labels are conserved and each dynamic block is
nonnegative. For $g\leq0$ the vacua have no quantum region; their
dynamic Hamiltonian is a sum of frozen-site penalties with positive
eigenvalues at least one. Every other wall block has offset at least
$r$, and a wrong tableau adds a unit penalty. Thus the gap is at least
$r$ on every torus.

Fix sufficiently small positive rational $g$, first off ties. Choose
$L=n!$ large enough for Corollary~\ref{cor:selection}. All nonminimizing
classical blocks then lie at least $5/8$ above the minimum. The finitely
specified word machine has written the required prefix in every cell
once $h$ exceeds its fixed timing cutoff. Each minimizing block has a
tensor sum of identical disjoint cell counters and frozen exterior
sites. Its first positive excitation is exactly $D(h)$. Since
$k(h)\geq3$, $D(h)\leq4\pi^2/64<5/8$, so no other classical sector
undercuts it. At a tie, minimizing stripe stacks use the two pitches;
the least positive dynamic excitation over all minimizing blocks is
$\min(D(p-1),D(p))$. Pure stacks show that both candidate values occur.
Rational wall energies and the finite timing cutoff give a computable
factorial index beyond which this identity holds, also at rational ties.

Every classical block has a zero-energy dynamic state, so the density
is the wall minimum. Equations~\eqref{eq:ruler-energy} and
\eqref{eq:ruler-asymptotic} prove its asserted expansion and
nonanalyticity at zero. The positive side has commensurate first-order
kinks accumulating at zero; the gap at zero itself remains positive.

For $g\in(g_p,g_{p-1})$,
$\log(1/g)=\log p-\log(2r)+o(1)$ and
$\ell(p-1)=\log_2p+O(1)$. Taking logarithms in
\eqref{eq:cell-bracket} gives
\[
 x(g)=2\beta\rho_{\ell(p_*(g)-1)}(A)
       +O\!\left(\frac{\log\log(1/g)}{\log(1/g)}\right).
\]
The pitch takes every sufficiently large integer. Its sampled bitlength
is nondecreasing, unbounded and advances by unit steps. Ties use the
minimum of two gaps; the two neighbouring normalized logarithms differ
by $o(1)$, so ties do not alter either endpoint. The endpoints are
therefore $2\beta$ times the two densities. Choose
$\beta\geq\max(1,b/2)$ and use Lemma~\ref{lem:word} with densities
$a/(2\beta),b/(2\beta)$. This proves attainment. The computably tamed
inner limit and Proposition~\ref{prop:critical-ceiling} give necessity.
\end{proof}

The canonical \emph{structural} length of this construction is its pitch:
$\xi(g)=p_*(g)=g^{-1+o(1)}$. If one defines a structural $\nu$ from this
length, it equals one, leaving the uncomputability in $z\nu$. Ground
space degeneracy means that correlation lengths of chosen ground states
need not equal this pitch. An equality with a GNS gap, or an assertion
about all torus sizes, requires a separate argument and is not used.
Different stripes may choose independent horizontal phases, so not
every minimizing configuration has one common vertical period.
The structural length denotes the cell pitch, or the period of
a chosen periodic representative.
